\chapter{الجسيمات المشحونة، وتحويلات الحقول، والموجات النسبية}

في المرة السابقة بدأنا نسأل كيف يتحرك جسيم مشحون في حقول كهربائية ومغناطيسية عندما تُفرَض مبادئ النسبية. وقانون قوة لورنتز المألوف يحتوي بالفعل على الدليل: فحده المغناطيسي يعتمد على سرعة الجسيم، ولذلك لا يمكن لمراقبين مختلفين أن يتفقوا، كلٌّ على حدة، على ما هو كهربائي وما هو مغناطيسي. ومهمتنا هي أن نصوغ القانون بحيث يستخدم كل مراقب عطالي المعادلة نفسها، رغم أن الحقول وسرعة الجسيم تتغير من إطار إلى آخر.

ما لم يُذكر خلاف ذلك، نستخدم وحدات يكون فيها \(c=1\) وتوقيع المترية \((-,-,-,+)\). وتُرتَّب مركبات الزمكان بحيث تأتي المركبات المكانية الثلاث أولًا:
\begin{equation}
  \mu=(1,2,3,0)=(x,y,z,t).
\end{equation}
أما المؤشران اللاتينيان \(m,n\) فلا يمتدان إلا على الاتجاهات المكانية الثلاثة.

\section{ما الذي يبقى من قانون نيوتن للقوة}

في الميكانيكا النيوتنية،
\begin{equation}
  \boldsymbol F=m\boldsymbol a
\end{equation}
له الصيغة نفسها في كل إطار عطالي لأن المراقبين الغاليليين يتفقون على التسارع. كما ينسبون إلى الجسيم الكتلة نفسها والقوة نفسها. وما إن تتغير الحركيات، فلا يوجد سبب يبرر توقع بقاء هذه العلاقة بعينها بين القوة والتسارع العادي من غير تغيير.

تنبهنا الكهرومغناطيسية والجاذبية بالفعل إلى أن التمديدات النسبية لا يلزم أن تشبه بعضها بعضًا. فقانون كولوم وقانون نيوتن للجاذبية كلاهما قانون تربيع عكسي، لكن الجاذبية لا يمكن إدماجها في النسبية الخاصة إدماجًا طبيعيًا من دون تغيير أعمق في هندسة الزمكان. أما الكهرومغناطيسية فيمكن ذلك. بل إن المغناطيسية نفسها أثر نسبي: فعندما نعيد إدخال عوامل \(c\)، تحمل القوة المغناطيسية عاملًا من السرعة مقسومًا على \(c\)، وتختفي في النهاية الصورية \(c\to\infty\).

الجزء الباقي من قانون نيوتن هو القول إن القوة هي معدل تغيّر الزخم. وعند ثبات كتلة السكون،
\begin{equation}
  \boldsymbol F
  =m\frac{d\boldsymbol v}{dt}
  =\frac{d}{dt}(m\boldsymbol v)
  =\frac{d\boldsymbol p}{dt}.
\end{equation}
ستغيّر النسبية العلاقة بين الزخم والسرعة، لكنها ستحفظ هذه الطريقة في تنظيم المعادلة.

\begin{figure}[H]
\centering
\includegraphics[width=0.64\textwidth]{lecture_07_figure_01.png}
\caption{إعادة كتابة القوة على صورة \(d\boldsymbol p/dt\) بينما يُكتَب الجزء الكهربائي من قانون لورنتز على اللوح.}
\lectureframe{00:08:20}
\end{figure}

في وحدات \(c=1\)، تكون الصيغة ما قبل النسبية لقانون لورنتز للقوة هي
\begin{equation}
  \frac{d\boldsymbol p}{dt}
  =q\bigl(\boldsymbol E+\boldsymbol v\times\boldsymbol B\bigr).
  \label{eq:lorentz-three}
\end{equation}
وقد يعتمد كلا المجالين على الموضع والزمن. والنهاية غير النسبية هي الآن النهاية \(v\ll 1\)، لا نهايةً يتغير فيها الثابت العددي \(c=1\) على نحو ما. والحد الكهربائي في \eqref{eq:lorentz-three} هو بالفعل متجه عادي. ولكي نفهم الحد المغناطيسي بصيغة يمكن نقلها إلى الزمكان، يجب أن ننظر من جديد إلى الضرب الاتجاهي.

\section{الضرب الاتجاهي بوصفه موترًا ضدّي التناظر}

خذ متجهين ثلاثيي الأبعاد \(\boldsymbol A\) و\(\boldsymbol B\). إن جداءهما الخارجي
\begin{equation}
  T_{mn}=A_mB_n
\end{equation}
هو موتر من الرتبة الثانية له تسع مركبات. ولا يملك، في العموم، أي تناظر. وأخذ الجزء ضدّي التناظر يعطي
\begin{equation}
  K_{mn}=A_mB_n-A_nB_m,
  \qquad K_{mn}=-K_{nm}.
  \label{eq:antisym-product}
\end{equation}
وتنعدم عناصره القطرية، بينما يتحدد كل عنصر تحت القطر بالعنصر المناظر فوقه. ولذلك لا يملك موتر \(3\times3\) ضدّي التناظر إلا ثلاث مركبات مستقلة.

هذا العدّ موحٍ، لكنه غير كافٍ بمفرده. فكون الكميات ثلاثًا لا يعني تلقائيًا أنها تتحول كمتجه. إلا أنه في ثلاثة أبعاد يمكن مطابقة العناصر المستقلة الثلاثة في \eqref{eq:antisym-product} مع مركبات الضرب الاتجاهي
\(\boldsymbol C=\boldsymbol A\times\boldsymbol B\). ويجري على اللوح أولًا وضع مؤقت لـ \(C_2\) في الموضع العلوي الأيمن:

\begin{figure}[H]
\centering
\includegraphics[width=0.75\textwidth]{lecture_07_figure_02.png}
\caption{الجداء ضدّي التناظر أثناء مطابقة عناصره المستقلة الثلاثة مع مركبات \(\boldsymbol C\). وما يزال الوسم العلوي الأيمن مؤقتًا هنا.}
\lectureframe{00:17:33}
\end{figure}

وبالتدوير \(1\to2\to3\to1\) نحصل على صيغ المركبات
\begin{align}
  C_3 &= A_1B_2-A_2B_1, \\
  C_1 &= A_2B_3-A_3B_2, \\
  C_2 &= A_3B_1-A_1B_3.
  \label{eq:cross-components}
\end{align}
ويُظهر السطر الأخير أن العنصر العلوي الأيمن هو \(-C_2\) لا \(C_2\). وهذا التصحيح جزء من الحجة، لا تفصيل ينبغي إخفاؤه.

\begin{figure}[H]
\centering
\includegraphics[width=0.62\textwidth]{lecture_07_figure_03.png}
\caption{الهويات الدورية المكتملة للمركبات التي تثبّت الإشارة المؤقتة في المصفوفة ضدية التناظر.}
\lectureframe{00:21:54}
\end{figure}

وتكون المصفوفة الناتجة هي
\begin{equation}
  K_{mn}=
  \begin{pmatrix}
    0 & C_3 & -C_2 \\
    -C_3 & 0 & C_1 \\
    C_2 & -C_1 & 0
  \end{pmatrix}.
  \label{eq:cross-matrix}
\end{equation}

توجد هندسة نافعة وراء هذا الترميز. فمركبة ذات المؤشرين \(1,2\) تنتمي بطبيعتها إلى مستوي \(xy\)، وتُطابَق مع مركبة متجهية على امتداد المحور \(z\) العمودي عليه. ويمكن وصف الشيء الموجَّه نفسه إما بمستوى وإما بسهم عمودي على ذلك المستوى. والوصف بالسهم خاص بثلاثة أبعاد؛ ففي الأبعاد الأعلى يبقى الموتر من الرتبة الثانية ضدّي التناظر موجودًا، لكنه لا يعود يملك عدد المركبات نفسه الذي يملكه المتجه. ومن ثم فإن الموتر، لا السهم العمودي، هو صورة الضرب الاتجاهي التي تتعمم.

وبصورة مكافئة، بعد اختيار اتجاهية،
\begin{equation}
  C_k=\frac12\epsilon_{kmn}K_{mn}.
\end{equation}
والحاجة إلى اختيار اتجاهية هي منشأ قاعدة اليد اليمنى، وكذلك التمييز المألوف بين المتجهات القطبية والمحورية.
\editorialnote{إن صيغة Levi--Civita ومصطلح المتجهات المحورية هما ترميزان قياسيان للهندسة التي تُقابِل بين المستوي والعمودي عليه، والتي طُوِّرت صراحةً في المحاضرة.}

\section{المجال المغناطيسي في الصيغة الموترية}

وينطبق الوصف المزدوج نفسه على المجال المغناطيسي. ونعرّف موتره المكاني ضدّي التناظر بـ
\begin{equation}
  B_{mn}=
  \begin{pmatrix}
    0 & -B_3 & B_2 \\
    B_3 & 0 & -B_1 \\
    -B_2 & B_1 & 0
  \end{pmatrix}.
  \label{eq:magnetic-tensor}
\end{equation}
والإشارة الكلية الإضافية بالنسبة إلى \eqref{eq:cross-matrix} هي اصطلاح. وبعد اختياره يجب استعماله باتساق؛ فتغيير إشارة المركبات المتجهية الثلاث جميعًا يصف المطابقة نفسها بين الموتر والمتجه، ولكن باصطلاح اتجاهية معاكس.

\begin{figure}[H]
\centering
\includegraphics[width=0.72\textwidth]{lecture_07_figure_04.png}
\caption{المجال المغناطيسي مكتوبًا بوصفه مصفوفة مكانية ضدية التناظر وكذلك بدلالة مركباته المتجهية الثلاث.}
\lectureframe{00:24:40}
\end{figure}

والآن خذ المركبة الثالثة من الضرب الاتجاهي:
\begin{align}
  (\boldsymbol v\times\boldsymbol B)_3
  &=v_1B_2-v_2B_1 \\
  &=v_1B_{13}+v_2B_{23}+v_3B_{33}.
\end{align}
ويمكن إضافة الحد الأخير لأن \(B_{33}=0\). وبتدوير المركبة الحرة نحصل على القاعدة الموجزة
\begin{equation}
  (\boldsymbol v\times\boldsymbol B)_m=v_nB_{nm},
  \label{eq:cross-contraction}
\end{equation}
حيث يُجمع على المؤشر المكرر \(n\). ومن ثم فإن القوة المغناطيسية هي
\begin{equation}
  F_m^{(B)}=qv_nB_{nm}.
\end{equation}

\begin{figure}[H]
\centering
\includegraphics[width=0.78\textwidth]{lecture_07_figure_05.png}
\caption{المركبة الثالثة من \(\boldsymbol v\times\boldsymbol B\) وتعبيرها العام بوصفها انكماشًا مع \(B_{nm}\).}
\lectureframe{00:28:30}
\end{figure}

\begin{classroomqa}
\audiencequestion{هل يوسم المؤشر الحر \(m\) مركبةً مغناطيسيةً منفصلة للقوة؟}
\lecturerresponse{لا. إنه يوسم أي مركبة مكانية من القوة يجري حسابها. وقد ضُبط المجال الكهربائي ببساطة على الصفر في الحجة الحالية، وكل مؤشر هنا لا يمتد إلا على الاتجاهات المكانية الثلاثة.}
\end{classroomqa}

يجعل الترميز الموترِي خاصية عدم إنجاز المجال المغناطيسي شغلًا شبه تلقائية. لنشتق مربع السرعة:
\begin{align}
  \frac{d}{dt}(v_mv_m)
  &=2\dot v_m v_m \\
  &\propto 2(\boldsymbol v\times\boldsymbol B)_m v_m \\
  &=2v_nB_{nm}v_m.
  \label{eq:speed-derivative}
\end{align}
فالحدود التي فيها \(m=n\) تنعدم لأن العناصر القطرية في \(B_{nm}\) تساوي الصفر. وكل حد خارج القطر له قرين تُبدَّل فيه \(m\) و\(n\)؛ ويتلاشى الاثنان لأن \(v_nv_m\) متناظر، بينما \(B_{nm}\) ضدّي التناظر.
وعليه
\begin{equation}
  \frac{d}{dt}(v_mv_m)=0.
\end{equation}
فقوة مغناطيسية صرفة تثني المسار من دون أن تغيّر سرعته أو طاقته الحركية.

\begin{figure}[H]
\centering
\includegraphics[width=0.68\textwidth]{lecture_07_figure_06.png}
\caption{الحدود المتزاوجة في \(v_nB_{nm}v_m\) تتلاشى لأن الموتر المغناطيسي يغيّر إشارته عندما يُبدَّل مؤشراه.}
\lectureframe{00:35:40}
\end{figure}

هذا البرهان أكثر إسهابًا عمدًا من مجرد القول إن الضرب الاتجاهي عمودي على عامله الأول. أما العبارة الأخيرة فخاصة بتمثيل المتجهات الثلاثي الأبعاد. وانكماش زوج متناظر مع موتر ضدّي التناظر هو الحجة التي ستبقى صالحة في أربعة أبعاد.

\begin{classroomqa}
\audiencequestion{هل يجب أن يكون المجال المغناطيسي مستقلًا عن الزمن لكي تبقى السرعة ثابتة؟}
\lecturerresponse{الجبر أعلاه لا يشتق المجال قط، ولا يتطلب \(\partial_t\boldsymbol B=0\). فالحد المغناطيسي وحده لا ينجز شغلًا. لكن في الكهروديناميكا، يولِّد المجال المغناطيسي المتغير مجالًا كهربائيًا، ويمكن لذلك المجال الكهربائي أن يغيّر سرعة الجسيم.}
\end{classroomqa}

\section{لماذا يجب أن يمتزج المجالان الكهربائي والمغناطيسي}

يمكن رؤية امتزاج \(\boldsymbol E\) و\(\boldsymbol B\) قبل أي حساب رباعي الأبعاد. لننظر في حركة بطيئة، بحيث تكون العبارات النيوتنية العادية عن التسارع كافية. ليكن المجال المغناطيسي مشيرًا على امتداد المحور \(y\)، وليتحرك جسيم مشحون على امتداد المحور \(x\). عندئذ تشير القوة المغناطيسية على امتداد المحور \(z\):
\begin{equation}
  \boldsymbol F=q\boldsymbol v\times\boldsymbol B.
\end{equation}

\begin{figure}[H]
\centering
\includegraphics[width=0.62\textwidth]{lecture_07_figure_07.png}
\caption{رسم اللوح لجسيم يتحرك عبر مجال مغناطيسي؛ وتُستخدَم الاتجاهات الثلاثة لمقارنة وصفي المختبر والإطار المصاحب للحركة.}
\lectureframe{00:42:20}
\end{figure}

يمكن إعادة رسم الهندسة من دون منظور قاعة الدرس:
\begin{figure}[H]
\centering
\begin{tikzpicture}[scale=1.05,>=Stealth]
  \draw[->] (0,0) -- (3.8,0) node[right] {\(x\)};
  \draw[->] (0,0) -- (-1.9,2.4) node[above left] {\(y\)};
  \draw[->] (0,0) -- (0,3.1) node[above] {\(z\)};
  \draw[->,very thick] (0,0) -- (2.7,0)
    node[midway,below] {\(\boldsymbol v\)};
  \draw[->,very thick] (0,0) -- (-1.35,1.7)
    node[midway,above left] {\(\boldsymbol B\)};
  \draw[->,very thick] (0,0) -- (0,2.35)
    node[midway,right] {\(\boldsymbol F,\ \boldsymbol E'\)};
\end{tikzpicture}
\caption{في المختبر تكون القوة الموجَّهة على امتداد \(z\) مغناطيسية؛ أما في الإطار الآني المصاحب للحركة، فلا بد أن تكون القوة نفسها كهربائية.}
\end{figure}

والآن أدخل راصدًا يتحرك مع الجسيم في تلك اللحظة. سرعة الجسيم في الإطار الجديد تساوي صفرًا، ولذلك يزول الحد المغناطيسي \(q\boldsymbol v'\times\boldsymbol B'\). ومع ذلك، في حد السرعات البطيئة لا بد أن يتفق الراصدان على أن للجسيم تسارعًا في اتجاه \(z\). وإذا كان قانون قوة لورنتز يحتفظ بالشكل نفسه في كلا الإطارين، فلا بد أن يعزو الراصد المصاحب القوة إلى مجال كهربائي \(\boldsymbol E'\) على امتداد \(z\). لذا فإن مجالًا كان مغناطيسيًا صرفًا في المختبر لا يكون مغناطيسيًا صرفًا في الإطار المتحرك.

هذا هو السبب الفيزيائي للبحث عن كائن واحد تشمل مركباته المجالين معًا. فتغيير الإطار لا يقتصر على تدوير \(\boldsymbol E\) بين مركباته الثلاث الخاصة به و\(\boldsymbol B\) بين مركباته الثلاث الخاصة به؛ بل يمزج المجموعتين.

\section{موتر المجال الكهرومغناطيسي}

يشكّل المجالان الكهربائي والمغناطيسي معًا ست مركبات. وموترٌ ضديد التناظر رباعي الأبعاد يملك كذلك ست مركبات مستقلة، لأن
\begin{equation}
  \frac{4(4-1)}{2}=6.
\end{equation}
وهذا لا يعني أن موترًا ضديد التناظر هو حرفيًا المجموعة الوحيدة من ست كميات يمكن اختراعها: فست كميات قياسية، أو متجه وكمّيتان قياسيتان، تحتوي أيضًا على ستة أعداد. إنه الكائن الهندسي الطبيعي الذي يمكن أن تمتزج مركباته على النحو المطلوب.

ولإبقاء اصطلاح اللوح صريحًا، نعرض الموتر المتباين بالترتيب المكاني أولًا \((1,2,3,0)\):
\begin{equation}
  F^{\mu\nu}=
  \begin{pmatrix}
    0 & -B_3 & B_2 & E_1 \\
    B_3 & 0 & -B_1 & E_2 \\
    -B_2 & B_1 & 0 & E_3 \\
    -E_1 & -E_2 & -E_3 & 0
  \end{pmatrix}.
  \label{eq:l07-field-tensor}
\end{equation}
والكتلة العلوية اليسرى \(3\times3\) هي بالضبط الموتر المغناطيسي \eqref{eq:magnetic-tensor}؛ أما المركبات التي لها دليل مكاني واحد ودليل زمني واحد فتحتوي على المجال الكهربائي. وتوفّر خاصية ضدّ التناظر باقي العناصر.

\begin{figure}[H]
\centering
\includegraphics[width=0.72\textwidth]{lecture_07_figure_08.png}
\caption{تجميع الكتلة المكانية المغناطيسية والعمود الزمني المكاني الكهربائي في موتر المجال رباعي الأبعاد.}
\lectureframe{00:46:35}
\end{figure}

\begin{classroomqa}
\audiencequestion{متى ينبغي أن تحمل مركبات المجال هذه مؤشرات علوية بدلًا من سفلية؟}
\lecturerresponse{بالنسبة إلى التسميات المكانية الثلاث وحدها، يهمل اللوح أحيانًا هذا التمييز عند الإشارة إلى حقول ماكسويل--فاراداي المألوفة. أمّا في حساب تغايري فلا بد من استعادة هذا التمييز: فرفع دليل مكاني أو خفضه يغيّر إشارة وفق اصطلاح المتريّة المختار.}
\end{classroomqa}

الثابت ليس ترتيبًا معيّنًا للإشارات على اللوح، بل موتر الزمكان الضديد التناظر مقترنًا باصطلاح متّسق للمتريّة والمؤشرات. وقد توحي الرمزية \(F\) بفاراداي، مع أن المحاضرة لا تجزم بهذا العزو التاريخي. وقد صاغ آينشتاين تحوّل المركبات الكهربائية والمغناطيسية في 1905؛ أما اللغة الموترة الرباعية الأبعاد المدمجة فجاءت مع الصياغة الزمكانية اللاحقة المرتبطة بمينكوفسكي.

\section{السرعة الرباعية والقوة الرباعية}

ويجب أيضًا جعل طرف المعادلة الخاص بالجسيم متغايرًا. فالزمن الخاص \(\tau\) ثابت، لذا فإن اشتقاق متجه الموضع الرباعي بالنسبة إليه يعطي متجهًا رباعيًّا آخر:
\begin{equation}
  u^\mu=\frac{dx^\mu}{d\tau}
  =\gamma(v_x,v_y,v_z,1),
  \qquad
  \gamma=\frac{1}{\sqrt{1-v^2}}.
  \label{eq:four-velocity}
\end{equation}
وتنتج المركبة الأخيرة من \(dt/d\tau=\gamma\). أما الزخم الرباعي فهو
\begin{equation}
  p^\mu=mu^\mu,
\end{equation}
وجزؤه المكاني هو الزخم النسبي
\(\boldsymbol p=m\gamma\boldsymbol v\).

\begin{figure}[H]
\centering
\includegraphics[width=0.75\textwidth]{lecture_07_figure_09.png}
\caption{السرعة الرباعية، ومركباتها، ومعامل لورنتز، والزخم الرباعي.}
\lectureframe{00:50:30}
\end{figure}

ولأن \(p^\mu\) متجه رباعي و\(\tau\) كمية قياسية، فإن الطرف الأيسر النسبي الطبيعي هو
\begin{equation}
  f^\mu\equiv\frac{dp^\mu}{d\tau}.
  \label{eq:four-force-def}
\end{equation}
ولا يمكن اختيار القوة الرباعية اعتباطًا. فالسرعة الرباعية لها معيار ثابت،
\begin{equation}
  u_\mu u^\mu=1.
\end{equation}
ويعطي الاشتقاق على طول خط العالم
\begin{align}
  0
  &=\frac{d}{d\tau}(u_\mu u^\mu) \\
  &=\frac{du_\mu}{d\tau}u^\mu
    +u_\mu\frac{du^\mu}{d\tau} \\
  &=2u_\mu\frac{du^\mu}{d\tau}.
\end{align}
ومع ثبات الكتلة يستلزم ذلك
\begin{equation}
  f^\mu u_\mu=0.
  \label{eq:force-orthogonal}
\end{equation}
وهكذا تكون القوة الرباعية متعامدة، بمعنى مينكوفسكي، مع السرعة الرباعية. وهذا قيد حركي بحت ناتج من تطبيع \(u^\mu\)، وليس خاصية مميزة للكهرومغناطيسية.

\begin{figure}[H]
\centering
\includegraphics[width=0.64\textwidth]{lecture_07_figure_10.png}
\caption{كون معيار السرعة الرباعية مساويًا للواحد، وما يترتب عليه من تعامد القوة الرباعية والسرعة الرباعية.}
\lectureframe{00:56:50}
\end{figure}

لقد زوّدنا الحساب المغناطيسي الثلاثي الأبعاد بالفعل بالحيلة اللازمة لتحقيق \eqref{eq:force-orthogonal}. لنُجرِ انكماشًا لموتر ضديد التناظر مع نسخة واحدة من السرعة:
\begin{equation}
  \boxed{
  \frac{dp^\mu}{d\tau}=qF^{\mu\nu}u_\nu.}
  \label{eq:covariant-lorentz}
\end{equation}
وفعلًا،
\begin{equation}
  f^\mu u_\mu
  =qF^{\mu\nu}u_\nu u_\mu=0,
\end{equation}
لأن \(F^{\mu\nu}\) ضديد التناظر و\(u_\nu u_\mu\) متناظر.
وهذا ليس قانون القوة النسبي الوحيد الممكن تصوّره، لكنه أبسط قانون خطي في كلٍّ من المجال الكهرومغناطيسي والسرعة.

\begin{figure}[H]
\centering
\includegraphics[width=0.72\textwidth]{lecture_07_figure_11.png}
\caption{معادلة قوة لورنتز المتغايرة على اليسار، وانكماشها المنعدم مع سرعة رباعية ثانية على اليمين.}
\lectureframe{01:00:00}
\end{figure}

\section{استعادة قوة لورنتز العادية}

يجب على معادلة تبدو متغايرة أن تعيد إنتاج قانون القوة المعروف. ومع اصطلاح المتريّة المذكور في البداية،
\begin{equation}
  u_\mu=\gamma(-v_x,-v_y,-v_z,1).
\end{equation}
فبالنسبة إلى المركبة \(x\)، تعطي \eqref{eq:l07-field-tensor} و\eqref{eq:covariant-lorentz}
\begin{align}
  \frac{dp^1}{d\tau}
  &=q\bigl(F^{12}u_2+F^{13}u_3+F^{10}u_0\bigr) \\
  &=q\gamma\bigl(B_3v_y-B_2v_z+E_1\bigr) \\
  &=q\gamma\bigl[\boldsymbol E+\boldsymbol v\times
    \boldsymbol B\bigr]_x.
  \label{eq:x-force}
\end{align}
وتتبع المركبات المكانية الأخرى بتدوير المؤشرات دوريًا.

\begin{classroomqa}
\audiencequestion{أين يقع الحد \(qE_x\) في التوسّع بالمركبات؟}
\lecturerresponse{إنه الحد \(qF^{10}u_0\). فمركبة موتر المجال ذات الدليل المكاني الواحد والدليل الزمني الواحد هي مركبة مجال كهربائي، و\(u_0=\gamma\).}
\end{classroomqa}

إن العامل \(\gamma\) في الطرف الأيمن ليس قوة إضافية. وبما أن
\begin{equation}
  \frac{dp^i}{d\tau}=\frac{dt}{d\tau}\frac{dp^i}{dt}
  =\gamma\frac{dp^i}{dt},
\end{equation}
فإنه يُختَزل من طرفي \eqref{eq:x-force}. ولذلك تكون المعادلة المكانية الدقيقة من جديد
\begin{equation}
  \frac{d\boldsymbol p}{dt}
  =q\bigl(\boldsymbol E+\boldsymbol v\times\boldsymbol B\bigr),
  \qquad
  \boldsymbol p=m\gamma\boldsymbol v.
  \label{eq:exact-three-force}
\end{equation}
ويبقى الشكل المألوف، لكن النسبية كامنة في الزخم.

\begin{figure}[H]
\centering
\includegraphics[width=0.75\textwidth]{lecture_07_figure_12.png}
\caption{قانون لورنتز ثلاثي الأبعاد بعد استبعاد عامل الزمن الخاص المشترك؛ ويحتفظ الزخم بعامله \(\gamma\).}
\lectureframe{01:06:30}
\end{figure}

وتحمل المركبة الزمنية العلاقة المناظرة الخاصة بالطاقة،
\begin{equation}
  \frac{dp^0}{dt}=q\boldsymbol E\cdot\boldsymbol v,
\end{equation}
ولذلك لا يبذل المجال المغناطيسي شغلاً، في حين قد يغيّر المجال الكهربائي طاقة الجسيم.
\editorialnote{التحقق من المركبة الزمنية هو النظير المباشر للحساب المكاني في المحاضرة؛ وقد أُدرج لإظهار خاصية انعدام الشغل نفسها في الصياغة الرباعية الأبعاد.}

\begin{classroomqa}
\audiencequestion{هل هذا الاشتقاق رياضيات أم فيزياء؟}
\lecturerresponse{إنه كلاهما. فجبر الموترات يزوّدنا بالشكل المدمج، لكن فيزياء جوهرية دخلت عبر المبدأ القائل إن جميع الأطر القصورية تخضع للقوانين نفسها، وعبر افتراض أن معادلات الحركة من الرتبة الثانية. فقانون من الرتبة الثانية يتنبأ بالمستقبل انطلاقًا من الموضع والسرعة؛ أما قانون من الرتبة الثالثة فيتطلب التسارع بوصفه معطيات ابتدائية إضافية. والتجربة، لا الترميز، هي التي تخبرنا أيَّ نوع من القوانين تستخدمه الطبيعة.}
\end{classroomqa}

\section{كيف تتحول المجالات}

للمعادلة \eqref{eq:covariant-lorentz} الصورة نفسها في كل إطار لأن طرفيها متجهان رباعيان. أما موتر المجال نفسه فيتحول بمصفوفة لورنتز واحدة لكل دليل:
\begin{equation}
  F'^{\mu\nu}
  =\Lambda^\mu{}_{\rho}\Lambda^\nu{}_{\sigma}F^{\rho\sigma}.
  \label{eq:tensor-transform}
\end{equation}
ويشتق اللوح هذه القاعدة بأن يبدأ أولًا بتحويل حاصل الضرب البسيط \(A^\mu B^\nu\). وبمجرد أن تُعرف القاعدة لحاصل ضرب متجهين، تعطي الخطية القاعدة لأي موتر من الرتبة الثانية.

بالنسبة إلى دفعة على طول المحور \(1\)، تتحول الإحداثيات وفقًا لـ
\begin{align}
  x'^1&=\gamma(x^1-vx^0), &
  x'^0&=\gamma(x^0-vx^1), \\
  x'^2&=x^2, & x'^3&=x^3.
\end{align}
طبّق التحويل نفسه على كلٍّ من دليلي الموتر على حدة. وبالنسبة إلى المركبة الموازية للدفعة، يبدأ اللوح بحاصل ضرب:
\begin{align}
  (A'^0B'^1)
  =\gamma^2\bigl(
     A^0B^1-vA^0B^0-vA^1B^1+v^2A^1B^0
   \bigr).
\end{align}
وباستخال كل حاصل ضرب بمركبة الموتر المناظرة ثم استخدام خاصية ضدّ التناظر، \(F^{00}=F^{11}=0\) و\(F^{10}=-F^{01}\)، نحصل على
\begin{equation}
  F'^{01}
  =\gamma^2(1-v^2)F^{01}
  =F^{01}.
\end{equation}
ومن ثم فإن مركبة المجال الكهربائي الموازية للحركة لا تتغير. وتتصرف المركبة المغناطيسية الموازية على نحو مماثل.

أما المركبات المستعرضة فمختلفة. ووفق اصطلاح الإطار ذي الشرطة والإشارات المستخدم على اللوح،
\begin{align}
  F'^{02}
  &=\gamma\bigl(F^{02}-vF^{12}\bigr), \\
  E'_2
  &=\gamma(E_2+vB_3).
  \label{eq:e-transform}
\end{align}
تكتسب مركبة كهربائية مستعرضة جزءًا من مركبة مغناطيسية. وتنتج قوانين التحويل المتبقية من قاعدة المؤشرين نفسها، وتُترك تمرينًا في المحاضرة.

\begin{figure}[H]
\centering
\includegraphics[width=0.76\textwidth]{lecture_07_figure_13.png}
\caption{تحويل الموتر المستعرض عند الخطوة التي يتبيّن فيها أن المركبة \(E'_2\) تحتوي على حد يتناسب مع \(vB_3\).}
\lectureframe{01:23:40}
\end{figure}

تعتمد الإشارة في \eqref{eq:e-transform} على أيِّ إطار يُعلَن أنه يتحرك بسرعة \(+v\)، وكذلك على اصطلاحي المؤشرات والمتريّة. لكن مضمونها الثابت لا يعتمد على ذلك: فالمركبات الكهربائية والمغناطيسية المستعرضة تمتزج، وتظل معادلة القوة المتغايرة صحيحة بعد التحويل.

\begin{classroomqa}
\audiencequestion{بالنسبة إلى تيار في سلك، هل يمكن التحرك بسرعة ما تجعل المجال المغناطيسي ينعدم؟}
\lecturerresponse{لا يكفي التحرك مع إلكترونات التوصيل، لأن شحنات الشبكة البلورية الموجبة تكون عندئذٍ متحركة في الإطار الجديد. ولذلك فالسلك المتعادل الحامل للتيار أدقّ من مجرد كتلة شحنة مفردة متحركة. ولا يمكن تحويل مجال مغناطيسي صرف بحيث يزول بالكامل. وإذا كان المجالان الكهربائي والمغناطيسي موجودين معًا، فقد تلغي أجزاؤهما المتحوّلة بعضها بعضًا في ظروف خاصة؛ أمّا في حالة شحنة واحدة تتحرك حركة منتظمة، فإن إطار سكونها هو الإطار البسيط الذي ينعدم فيه مجالها المغناطيسي.}
\end{classroomqa}

يعتمد الاختبار الدقيق القائم على الكميات غير المتغيرة على
\begin{equation}
  \boldsymbol B^2-\boldsymbol E^2,
  \qquad
  \boldsymbol E\cdot\boldsymbol B.
\end{equation}
ولا يمكن أن يوجد إطار له \(\boldsymbol B'=0\) إلا لحقل مهيمن كهربائيًا يحقق \(\boldsymbol E\cdot\boldsymbol B=0\)؛ أما الحقل الذي يكون مغناطيسيًا صرفًا في أحد الأطر فهو مهيمن مغناطيسيًا ولا يمكنه استيفاء ذلك الشرط.
\editorialnote{تبيّن هذه المعايير غير المتغيرة على نحو دقيق القيد الذي توصّلنا إليه في مناقشة الصف؛ ويُؤجَّل اشتقاقها.}

\begin{classroomqa}
\audiencequestion{هل يشكّل الكمون القياسي الكهربائي والكمون المتجهي المغناطيسي متجهًا رباعيًّا؟}
\lecturerresponse{نعم. ومع إشارات تعتمد على الاصطلاح، يندمجان في الكمون الرباعي \(A^\mu\). وستستخدم معادلات ماكسويل ذلك المقدار على نحو منهجي.}
\end{classroomqa}

\section{مشتقات الحقول}

نترك الآن مسار جسيم واحد وننظر في حقل عبر الزمكان كله. ولتكن
\begin{equation}
  \phi=\phi(x^1,x^2,x^3,x^0)
\end{equation}
حقلًا قياسيًا. وتشكّل مشتقاته الأربعة المركبات الموافقة للتغاير لمتجه رباعي:
\begin{equation}
  \partial_\mu\phi
  \equiv\frac{\partial\phi}{\partial x^\mu}
  =\left(
    \frac{\partial\phi}{\partial x},
    \frac{\partial\phi}{\partial y},
    \frac{\partial\phi}{\partial z},
    \frac{\partial\phi}{\partial t}
  \right).
\end{equation}
وتبرهن قاعدة السلسلة، عند تطبيقها على تحويل إحداثيات لورنتز، قانون التحول الموافق للتغاير هذا. أما رفع الفهرس فيعكس الإشارات المكانية لكنه يترك إشارة الزمن من دون تغيير:
\begin{equation}
  \partial^\mu\phi
  =\left(
    -\frac{\partial\phi}{\partial x},
    -\frac{\partial\phi}{\partial y},
    -\frac{\partial\phi}{\partial z},
    \frac{\partial\phi}{\partial t}
  \right).
\end{equation}

كل اشتقاق يضيف فهرسًا سفليًا. ومن ثم فإن \(\partial_\nu V^\mu\) موتر من الرتبة الثانية له فهرس علوي واحد وآخر سفلي، كما أن اشتقاق موتر ينتج مقدارًا له فهرس سفلي إضافي واحد. وهذه هي القاعدة الأساسية اللازمة لكتابة معادلات الحقول من دون اختيار إطار قصوري مفضّل.

\section{المعادلة الموجية القياسية غير المتغيرة}

أبسط المعادلات الموجية تحتوي على مشتقتين. وإذا أريد للمعادلة أن يكون لها الشكل نفسه في كل إطار قصوري، فينبغي أن يكون طرفها الأيسر مقدارًا قياسيًا. ويعطي عقد فهارس المشتقات الاختيار الأبسط الوحيد:
\begin{equation}
  \boxed{\partial^\mu\partial_\mu\phi=0.}
  \label{eq:wave-covariant}
\end{equation}
وبكتابة المركبات صراحةً وفق التوقيع المعتمد هنا،
\begin{equation}
  \frac{\partial^2\phi}{\partial t^2}
  -\frac{\partial^2\phi}{\partial x^2}
  -\frac{\partial^2\phi}{\partial y^2}
  -\frac{\partial^2\phi}{\partial z^2}=0.
  \label{eq:wave-components}
\end{equation}
فإن العقد في \eqref{eq:wave-covariant} يُظهر ثبات لورنتز بوضوح؛ والمعادلة \eqref{eq:wave-components} هي العبارة نفسها مكتوبةً مركبةً بمركبة.

\begin{figure}[H]
\centering
\includegraphics[width=0.75\textwidth]{lecture_07_figure_14.png}
\caption{الصيغة ذات عقد المشتقات للمعادلة الموجية القياسية، وتوسيعها على المركبات بصيغة الزمن ناقص المكان.}
\lectureframe{01:39:40}
\end{figure}

\begin{classroomqa}
\audiencequestion{هل يجب أن يكون الطرف الأيمن للمعادلة الموجية صفرًا؟}
\lecturerresponse{لا. يمكن أن يكون أي مقدار قياسي متوافق مع المسألة. فضرب ثابت في \(\phi\)، أو دالة قياسية لاخطية مثل دالة تتناسب مع \(\phi^2\)، يظل يحافظ على التغاير. ويعطي الصفر أبسط معادلة موجية حرة.}
\end{classroomqa}

الغاية من هذا المثال القياسي هي التمهيد. فمعادلات ماكسويل تمزج المجالين الكهربائي والمغناطيسي على نحو أشد تعقيدًا، لكنها هي أيضًا يمكن أن تُكتب بحيث يظهر ثبات لورنتز من بنيتها الفهرسية. وتلك هي المسألة التي سنعود إليها في المرة القادمة.
